% Bewertungspaket zum Gesamttest «Einheitskreis» — Quelle des PDFs (python3 scripts/build-lp-pdf.py einheitskreis).
% Ganze Punkte je Teilschritt; Werte mit python3 nachgerechnet (scripts/lp/einheitskreis/zahlen.py, 07.10.2026, Prüfung 08.10.2026,
% Folgefehler-Fälle eingeschlossen).
\documentclass[11pt]{article}
\usepackage{../lp-druck}
\renewcommand{\lpname}{Einheitskreis}
\renewcommand{\lpteilgebiet}{5.4}
\renewcommand{\lpfuss}{Bewertungspaket}
\newcolumntype{P}{>{\centering\arraybackslash}p{10mm}}
\newenvironment{raster}{\par\smallskip\noindent\tabularx{\linewidth}{@{}>{\raggedright\arraybackslash}p{38mm}P X@{}}\toprule
  \small\textsc{Teilschritt} & \small\textsc{P} & \small\textsc{Musterlösung}\\\midrule}{\bottomrule\endtabularx\par}
\newcommand{\fehler}[1]{\par\smallskip{\small\textcolor{lprot}{\bfseries Typische Fehler:} #1}}
\newcommand{\E}{\,{\footnotesize\bfseries\color{lpgruen}(E)}}
\newcommand{\A}{\,{\footnotesize\bfseries\color{lpblau}(A)}}
\begin{document}
\lpkopf{Bewertungspaket zum Gesamttest}{}

\begin{lpachtung}{Erst lösen, dann öffnen}
Dieses Paket enthält die Musterlösung. Wer es vor dem Lösen liest, misst am Ende nur, wie gut das Abschreiben gelingt.
\end{lpachtung}

\section*{1 · So gehst du vor}
\begin{enumerate}[leftmargin=*,itemsep=2pt]
\item \textbf{Gesamttest lösen} — vollständig, mit Weg und Skizzen, ohne dieses Paket; Teil A ohne, Teil B mit Taschenrechner.
\item \textbf{Fotografieren oder scannen}, gut lesbar, eine Seite pro Bild. \textbf{Kein Name} auf den Bildern.
  Handyfotos tragen oft unsichtbar Standort, Zeit und Gerät mit (Metadaten). Lade darum lieber einen Scan oder ein
  Bildschirmfoto des Fotos hoch, oder entferne vorher die Standortangaben.
\item \textbf{Bewerten lassen.} Ob du dafür eine KI nutzt und welche, entscheidest du selbst und in eigener Verantwortung —
  je nachdem, was dir aufgrund deines Alters und deiner persönlichen Situation erlaubt ist (Altersgrenzen und
  Nutzungsbedingungen des Anbieters, allenfalls Einverständnis deiner Eltern). Lade nur deine Lösung und dieses PDF hoch,
  keine persönlichen Daten, und füge den Auftrag aus Abschnitt~2 ein. Ohne KI kannst du dich mit Abschnitt~3 selbst bewerten.
\item \textbf{Rückmeldung prüfen.} Stimmt, was die KI aus deiner Handschrift und deinen Skizzen gelesen hat? Eine KI kann sich
  irren — im Zweifel gilt das Raster in Abschnitt~3.
\item \textbf{Gezielt wiederholen} nach Abschnitt~4.
\end{enumerate}

\needspace{8\baselineskip}
\section*{2 · Auftrag an die KI}
\begin{tcolorbox}[breakable,colback=lplinie!25,colframe=lplinie,boxrule=0.5pt,arc=1mm,fontupper=\small]
Du bist Korrektorin oder Korrektor für Mathematik an einer Schweizer Berufsmaturitätsschule. Im Anhang findest du (1)~das Bewertungspaket zum Gesamttest «Einheitskreis» mit Musterlösung und Punkteraster und (2)~Fotos meiner handschriftlichen Lösung.\par\medskip
Geh so vor:\par
1. Schreib zu jeder Aufgabe G1 bis G8 zuerst ab, was du in meiner Lösung liest — Weg, Ergebnis und was in den Skizzen eingezeichnet ist (wo liegt der Punkt, welche Strecken sind markiert). Was du nicht sicher lesen kannst, markierst du mit [unsicher]. Nicht raten.\par
2. Vergib die Punkte genau nach dem Raster im Bewertungspaket (Abschnitt 3), ganze Punkte je Zeile. Ziehe einen Fehler nur einmal ab: Rechne ich mit einem falschen Zwischenergebnis richtig weiter, gibt es die Folgepunkte. Ausnahme: Zeilen mit \textbf{(E)} gibt es nur für das richtige Ergebnis, nie als Folgepunkt. Die Zeilen «Typische Fehler» sagen, welche Rasterzeile bei diesem Fehler 0 Punkte gibt — sie sind kein zusätzlicher Abzug.\par
3. Andere richtige Wege zählen voll. Gleichwertige Schreibweisen zählen voll: Dezimalpunkt oder Dezimalkomma, $\frac{\sqrt3}{3}$ oder $\frac{1}{\sqrt3}$, $\frac{\sqrt2}{2}$ oder $\frac{1}{\sqrt2}$, Bruch oder exakte Dezimalzahl ($\frac{1}{2} = 0.5$), Winkel in Grad oder im Bogenmass — ausser wo das Raster eine bestimmte Form verlangt. In Teil A gilt ein gerundeter Dezimalwert ($0.866$ statt $\frac{\sqrt3}{2}$) nicht als exakt. Drei Stufen: Zeilen mit \textbf{(A)} verlangen nur Ablesen, diese Punkte gibt es auch ohne Weg. Zeilen mit \textbf{(E)} sind Ergebnispunkte: Das Ergebnis muss stimmen \emph{und} aus einem erkennbaren Weg hervorgehen, ausser die Zeile sagt ausdrücklich etwas anderes. Alle übrigen Zeilen bewerten den Weg selbst.\par
4. Begründe jeden Abzug in einem Satz und nenne die Fehlerart (zum Beispiel «Vorzeichen des Quadranten übersehen» oder «Sinus und Cosinus vertauscht»). Schreib die Musterlösung nicht ab — zeig mir, wo mein Fehler liegt.\par
5. Gib zum Schluss eine Tabelle «Aufgabe | Punkte | Begründung», die Gesamtpunktzahl (von 25) und — nach der Selbsteinschätzung in Abschnitt 4 — welches Kapitel des Leitprogramms ich wiederholen soll. Keine Note.\par
6. Fehlt ein Foto oder ist eine Aufgabe unlesbar, sag es, statt sie zu bewerten.
\end{tcolorbox}

\needspace{8\baselineskip}
\section*{3 · Musterlösung und Punkteraster}
\textbf{Allgemein:} Ganze Punkte je Zeile. Ein Fehler wird nur einmal abgezogen (Folgefehler). Gleichwertige Wege und
Schreibweisen zählen voll. In der Spalte P:
\textbf{(A)}~= nur ablesen, zählt auch ohne Weg · \textbf{(E)}~= Ergebnispunkt, Ergebnis richtig
\emph{und} Weg erkennbar (ausser die Zeile sagt ausdrücklich etwas anderes) · ohne Zeichen = die Zeile bewertet den Weg.
Teil A ist ohne, Teil B mit Taschenrechner gestellt; Lineal und Geodreieck sind in beiden Teilen erlaubt. Als Weg zählt auch eine Skizze am Einheitskreis, die den Quadranten,
den Referenzwinkel oder die Spiegelung zeigt.
«Typische Fehler» nennt, welche Zeile 0 Punkte gibt, und die Punkte, die bleiben.
\textbf{Teilrichtige Zeilen:} Verlangt eine Zeile mehrere Angaben und stimmt nur ein Teil, gibt die Zeile 0 Punkte —
ausser die Zeile nennt ausdrücklich eine Aufteilung.

\begin{lpaufgabe}{G1}{Am Einheitskreis zeichnen}{3 P}
\begin{raster}
(a) Punkt und Strecken & 1 & $P$ auf dem Kreis im III.~Quadranten, $40^\circ$ unter der negativen $x$-Achse: $P \approx (-0.77 \mid -0.64)$. Toleranz: $\pm 5^\circ$ (Geodreieck erlaubt). Die Strecke von $O$ nach $(\cos\varphi \mid 0)$ auf der $x$-Achse ist $\cos\varphi$, die senkrechte Strecke von dort zu $P$ ist $\sin\varphi$; beide beschriftet oder farbig unterschieden.\\
(b) Punkt $S$ & 1 & Die Gerade durch $P$ und $O$, über $O$ hinaus verlängert, trifft die Tangente \emph{über} der $x$-Achse. Zum richtigen $P$ ist $S \approx (1 \mid 0.84)$. Bewertet wird, ob $S$ auf der Tangente und auf der Verlängerung der \emph{eigenen} Geraden durch $P$ und $O$ liegt (Toleranz ein halbes Kästchen) — auch bei ungenauem $P$ (Folgepunkt).\\
(c) Vorzeichen & 1\A & $\sin 220^\circ < 0$, $\cos 220^\circ < 0$, $\tan 220^\circ > 0$. Alle drei nötig.\\
\end{raster}
\fehler{$P$ bei $140^\circ$ oder $320^\circ$ (falscher Quadrant): Zeile (a) 0; (b) und (c) gelten als Folgefehler, wenn sie zu diesem Punkt passen.
Sinus und Cosinus vertauscht beschriftet: Zeile (a) 0.
$S$ unter der $x$-Achse ($\approx (1 \mid -0.84)$, Strahl von $O$ in die falsche Richtung): Zeile (b) 0.
(c) $\tan 220^\circ < 0$: Zeile (c) 0.}
\end{lpaufgabe}

\begin{lpaufgabe}{G2}{Besondere Winkel}{4 P}
\begin{raster}
(a) & 1\E & $-135^\circ$ dreht im Uhrzeigersinn, derselbe Punkt wie $225^\circ$: III.~Quadrant, Referenzwinkel $45^\circ$: $\sin(-135^\circ) = -\frac{\sqrt2}{2}$\\
(b) & 1\E & $690^\circ - 360^\circ = 330^\circ$: IV.~Quadrant, Referenzwinkel $30^\circ$: $\cos 690^\circ = \frac{\sqrt3}{2}$\\
(c) & 1\E & $120^\circ$: II.~Quadrant, Referenzwinkel $60^\circ$: $\tan 120^\circ = \frac{\sqrt3/2}{-1/2} = -\sqrt3$ (im II.~Quadranten negativ)\\
(d) & 1\E & $-\frac{2\pi}{3} = -120^\circ$, derselbe Punkt wie $240^\circ$: III.~Quadrant, Referenzwinkel $60^\circ$: $\cos\left(-\frac{2\pi}{3}\right) = -\frac12$\\
\end{raster}
\fehler{Vorzeichen des Quadranten übersehen ($\frac{\sqrt2}{2}$ in (a), $-\frac{\sqrt3}{2}$ in (b), $\sqrt3$ in (c), $\frac12$ in (d)): die betroffene Zeile 0.
Referenzwinkel zur $y$-Achse gemessen, darum Sinus und Cosinus vertauscht ($\frac12$ in (b), $-\frac{\sqrt3}{3}$ in (c), $-\frac{\sqrt3}{2}$ in (d)): die Zeile 0.
(c) $-\frac{\sqrt3}{3}$ (Cosinus durch Sinus): Zeile (c) 0.
(d) über $\cos(-\alpha) = \cos\alpha$ zu $\cos 120^\circ = -\frac12$: ein richtiger Weg, Zeile (d) voll.
Gerundete Dezimalzahl statt exaktem Wert: die Zeile 0 (Teil A verlangt exakte Werte).
Als (E)-Weg genügen Quadrant und Referenzwinkel oder eine Skizze; fehlen beide, gibt die Zeile 0.}
\end{lpaufgabe}

\begin{lpaufgabe}{G3}{Trigonometrischer Pythagoras}{3 P}
\begin{raster}
Quadrant & 1 & $\sin\varphi > 0$: I.~oder II.~Quadrant; $\tan\varphi < 0$: II.~oder IV.~Quadrant. Also II.~Quadrant, $\cos\varphi < 0$. Quadrant \emph{und} Begründung aus beiden Vorzeichen nötig.\\
$\cos\varphi$ & 1\E & $\cos^2\varphi = 1 - \sin^2\varphi = 1 - \frac{64}{289} = \frac{225}{289}$, $\cos\varphi = \pm\frac{15}{17}$; im II.~Quadranten negativ: $\cos\varphi = -\frac{15}{17}$\\
$\tan\varphi$ & 1 & $\tan\varphi = \frac{\sin\varphi}{\cos\varphi} = \frac{8/17}{-15/17} = -\frac{8}{15}$. Folgepunkt: gilt auch mit einem falschen $\cos\varphi$, wenn der Quotient richtig gebildet ist.\\
\end{raster}
\fehler{I.~Quadrant ($\tan\varphi < 0$ übersehen), darum $\cos\varphi = \frac{15}{17}$: Zeilen Quadrant und $\cos\varphi$ 0, $\tan\varphi = \frac{8}{15}$ als Folgefehler 1\,P $\to$ 1 von 3\,P.
Quadrant richtig, aber $\cos\varphi = \frac{15}{17}$ (Vorzeichen nach dem Wurzelziehen vergessen): Zeile $\cos\varphi$ 0, $\tan\varphi = \frac{8}{15}$ als Folgefehler 1\,P $\to$ 2 von 3\,P.
$\cos\varphi = -\left(1 - \frac{8}{17}\right) = -\frac{9}{17}$ (ohne Quadrate): Zeile $\cos\varphi$ 0, $\tan\varphi = -\frac{8}{9}$ als Folgefehler 1\,P.
$\tan\varphi = -\frac{15}{8}$ (Cosinus durch Sinus): Zeile $\tan\varphi$ 0.
Gerundete Dezimalzahlen ($-0.882$, $-0.533$) statt exakter Werte: die betroffene Zeile 0 (Teil A, exakt verlangt).}
\end{lpaufgabe}

\begin{lpaufgabe}{G4}{Symmetrien}{4 P}
\begin{raster}
(a) & 1\E & $\sin 328^\circ = \sin(360^\circ - 32^\circ) = -\sin 32^\circ \approx -0.530$ — Spiegelung an der $x$-Achse\\
(b) & 1\E & $\cos 148^\circ = \cos(180^\circ - 32^\circ) = -\cos 32^\circ \approx -0.848$ — Spiegelung an der $y$-Achse (Supplement)\\
(c) & 1\E & $\cos 58^\circ = \cos(90^\circ - 32^\circ) = \sin 32^\circ \approx 0.530$ — Spiegelung an der Geraden $y = x$ (Komplement)\\
(d) & 1 & Falsch ist der letzte Schritt: Die Spiegelung am Ursprung ($180^\circ + \alpha$) kippt beide Koordinaten, also auch den Cosinus. Richtig: $\cos 212^\circ = -\cos 32^\circ \approx -0.848$. Fehler benannt \emph{und} richtiger Wert nötig.\\
\end{raster}
\textbf{Symmetrie genannt} heisst: die Spiegelung (an der $x$-Achse, an der $y$-Achse, an der Geraden $y = x$, am Ursprung) \emph{oder} die Formel (z.\,B. $\sin(360^\circ - \alpha) = -\sin\alpha$) \emph{oder} eine Skizze mit beiden Punkten. Eines davon genügt.
\fehler{(a) $0.530$ (Vorzeichen nicht gekippt): Zeile (a) 0.
(b) $0.848$ (Vorzeichen vergessen) oder $0.530$ (Sinus statt Cosinus): Zeile (b) 0.
(c) $0.848$ (Komplement und Supplement verwechselt, die Funktion nicht getauscht): Zeile (c) 0.
(d) nur «stimmt nicht» ohne Grund oder nur $-0.848$ ohne Fehler: Zeile (d) 0.
Richtige Zahl in (a) bis (c) ohne genannte Symmetrie und ohne Skizze: die Zeile 0 — (E) verlangt einen erkennbaren Weg.
Mit dem Taschenrechner gerechnet ($\sin 328^\circ \approx -0.5299$): Teil A ist ohne Rechner gestellt; gilt nur mit genannter Symmetrie.}
\end{lpaufgabe}

\begin{lpaufgabe}{G5}{Am Einheitskreis begründen}{2 P}
\begin{raster}
(a) gleicher Tangens & 1 & Der Punkt zu $\varphi + 180^\circ$ liegt dem Punkt zu $\varphi$ gegenüber (Spiegelung am Ursprung), auf derselben Geraden durch $O$. Diese Gerade trifft die Tangente $x = 1$ im selben Punkt $S$. \emph{Oder:} Beide Koordinaten kippen das Vorzeichen, der Quotient $\frac{\sin}{\cos}$ bleibt gleich.\\
(b) kein Tangens & 1 & Bei $90^\circ$ und $270^\circ$: $P$ liegt auf der $y$-Achse, die Gerade durch $O$ und $P$ ist parallel zur Tangente $x = 1$ und trifft sie nicht. \emph{Oder:} Dort ist $\cos\varphi = 0$, und durch null kann man nicht teilen. Beide Winkel \emph{und} ein Grund nötig.\\
\end{raster}
\fehler{(a) nur ein Zahlenbeispiel ($\tan 30^\circ = \tan 210^\circ$): Zeile (a) 0 — ein Beispiel begründet keine allgemeine Aussage.
(b) nur $90^\circ$: Zeile (b) 0. $0^\circ$ und $180^\circ$ (wo der Sinus null ist; dort ist $\tan\varphi = 0$): Zeile (b) 0.
(b) «$\tan 90^\circ = 0$»: Zeile (b) 0.}
\end{lpaufgabe}

\begin{lpaufgabe}{G6}{Das Schiff um die Boje}{3 P}
\begin{raster}
(a) Ost-Versatz & 1\E & $\cos 235^\circ \approx -0.574$~SM (also $0.574$~SM westlich)\\
(a) Nord-Versatz & 1\E & $\sin 235^\circ \approx -0.819$~SM (also $0.819$~SM südlich)\\
(b) Richtung & 1\A & südwestlich der Boje. Als Folgepunkt: die Richtung, die zu den eigenen Vorzeichen in (a) passt.\\
\end{raster}
\fehler{Ost- und Nord-Versatz vertauscht ($-0.819$ östlich, $-0.574$ nördlich): beide (a)-Zeilen 0, (b) als Folgefehler 1\,P $\to$ 1 von 3\,P.
Rechner im Bogenmass ($\cos 235 \approx -0.814$, $\sin 235 \approx 0.581$): beide (a)-Zeilen 0; (b) gilt als Folgefehler, wenn die Richtung zu den Vorzeichen passt.
Abweichung in der dritten Dezimale durch Runden zählt voll.}
\end{lpaufgabe}

\begin{lpaufgabe}{G7}{Was der Rechner liefert}{3 P}
\begin{raster}
(a) Punkt & 1 & $P$ im IV.~Quadranten, $63.4^\circ$ im Uhrzeigersinn unter der positiven $x$-Achse: $P \approx (0.45 \mid -0.89)$. Toleranz: ein Kästchen.\\
(b) zweiter Punkt & 1 & Der Punkt gegenüber, am Ursprung gespiegelt (II.~Quadrant, $116.6^\circ$), liegt auf derselben Geraden durch $O$ und hat denselben Tangens. Der Rechner liefert nur Winkel aus $]{-90^\circ};\, 90^\circ[$, der rechten Kreishälfte; der zweite Punkt liegt links. Lage (eingezeichnet oder beschrieben) \emph{und} Grund nötig, ein Winkel für den zweiten Punkt nicht.\\
(c) Winkel & 1\E & $-63.4^\circ + 360^\circ = 296.6^\circ$\\
\end{raster}
\fehler{(a) $P$ bei $63.4^\circ$ über der $x$-Achse (Vorzeichen übersehen): Zeile (a) 0.
(b) $180^\circ - (-63.4^\circ) = 243.4^\circ$ (Regel des Sinus; dort ist $\tan 243.4^\circ \approx 2$) oder ein anderer Punkt mit gleicher Höhe (Sinus-Symmetrie statt Tangens): Zeile (b) 0.
(c) $63.4^\circ$ oder $243.4^\circ$: Zeile (c) 0; $116.6^\circ$ ist der zweite Punkt, nicht $P$: Zeile (c) 0.}
\end{lpaufgabe}

\begin{lpaufgabe}{G8}{Nur ein Winkel?}{3 P}
\begin{raster}
(a) Skizze & 1 & Senkrechte $x = 0.7$; sie trifft den Kreis in zwei Punkten, symmetrisch zur $x$-Achse: bei $45.6^\circ$ und bei $-45.6^\circ$ ($= 314.4^\circ$).\\
(a) Begründung & 1 & «Nein» und ein Grund: Es gibt einen zweiten Kreispunkt mit der $x$-Koordinate $0.7$ (unter der $x$-Achse), \emph{oder}: Nach jeder vollen Runde ($360^\circ$) kommt man wieder zu denselben Punkten. Einer der beiden Gründe genügt; der Hinweis, dass der Rechner nur den Hauptwert aus $[0^\circ;\, 180^\circ]$ liefert, ist richtig, aber nicht nötig.\\
(b) Winkel & 1\E & z.\,B. $-45.6^\circ$ (Spiegelung an der $x$-Achse, $\cos(-\alpha) = \cos\alpha$) und $45.6^\circ + 360^\circ = 405.6^\circ$. Andere richtige Winkel zählen auch (etwa $-314.4^\circ$ oder $765.6^\circ$; Probe: der Rechner gibt $\cos \approx 0.7$). Ein negativer Winkel \emph{und} einer über $360^\circ$ nötig.\\
\end{raster}
\fehler{Waagrechte $y = 0.7$ gezeichnet (Sinus statt Cosinus): Skizzen-Zeile 0; eine dazu passende Begründung zählt.
«Ja, der Rechner hat es ausgerechnet»: Begründungs-Zeile 0.
(b) $-45.6^\circ$ fehlt und stattdessen $180^\circ - 45.6^\circ = 134.4^\circ$ (Regel des Sinus; $\cos 134.4^\circ \approx -0.7$): Zeile (b) 0.}
\end{lpaufgabe}

\needspace{14\baselineskip}
\section*{4 · Selbsteinschätzung}
\begin{tabularx}{\linewidth}{@{}l X@{}}\toprule
22 – 25 P & Die geprüften Teile sitzen. Wo du Punkte verloren hast: das Kapitel dieser Aufgabe nochmals (Zuordnung unten).\\
17 – 21 P & Den schwächsten Teil nochmals: Tüfteln und Übungen des Kapitels, in dem du die meisten Punkte verloren hast.\\
11 – 16 P & Zurück zu den Kapiteln aller Aufgaben, in denen du Punkte verloren hast.\\
0 – 10 P & Zurück zu Kapitel 1 und von dort der Reihe nach weiter.\\\midrule
\multicolumn{2}{@{}p{\linewidth}@{}}{\small\textbf{Aufgabe $\to$ Kapitel:} G1 $\to$ Kapitel 1, 3; G2 $\to$ Kapitel 2, 3, 5; G3 $\to$ Kapitel 1, 3; G4 $\to$ Kapitel 4; G5 $\to$ Kapitel 3, 4, 5; G6 $\to$ Kapitel 1; G7 $\to$ Kapitel 5; G8 $\to$ Kapitel 4, 5}\\\bottomrule
\end{tabularx}
\end{document}
